\documentclass[a4paper,12pt]{article}

\usepackage[T2A]{fontenc}
\usepackage[utf8]{inputenc}
\usepackage[russian]{babel}

\usepackage{amsmath,amssymb,mathtools}

\usepackage{helvet}
\renewcommand{\familydefault}{\sfdefault}

\usepackage{icomma}
\usepackage[a4paper,margin=2.2cm,headheight=16pt]{geometry}
\usepackage{microtype}

\usepackage{tikz}
\usetikzlibrary{arrows.meta,positioning,calc}

\usepackage{tabularx,booktabs}
\usepackage{enumitem}
\usepackage[hidelinks,unicode]{hyperref}
\usepackage{parskip}
\usepackage{fancyhdr}
\usepackage{setspace}
\usepackage{xcolor}

\definecolor{accent}{HTML}{008080}
\definecolor{darkaccent}{HTML}{004D4D}
\definecolor{textgray}{HTML}{333333}
\definecolor{softgray}{HTML}{EEF4F4}
\definecolor{muted}{HTML}{6D7B7B}

\onehalfspacing
\setlength{\parindent}{0pt}
\setlength{\parskip}{0.55em}
\setlength{\emergencystretch}{2em}

\setlist[itemize]{
    leftmargin=1.7em,
    itemsep=0.25em,
    topsep=0.3em
}

\setlist[enumerate]{
    leftmargin=2em,
    itemsep=0.55em,
    topsep=0.4em
}

\pagestyle{fancy}
\fancyhf{}
\fancyhead[L]{\small\textcolor{muted}{Кирилл Зиновьев}}
\fancyhead[R]{\small\textcolor{muted}{Проценты}}
\fancyfoot[C]{\small\thepage}
\renewcommand{\headrulewidth}{0.4pt}
\renewcommand{\headrule}{
    \hbox to\headwidth{\color{accent}\leaders\hrule height \headrulewidth\hfill}
}

\newcommand{\sectiontitle}[1]{%
    \vspace{0.6em}
    {\Large\bfseries\color{darkaccent}#1\par}
    \vspace{0.15em}
    {\color{accent}\rule{\linewidth}{0.6pt}}
    \vspace{0.4em}
}

\newcommand{\subsectiontitle}[1]{%
    \vspace{0.35em}
    {\large\bfseries\color{accent}#1\par}
    \vspace{0.15em}
}

\newcommand{\key}[1]{\textbf{\textcolor{darkaccent}{#1}}}

\begin{document}

% ============================================================
% ТИТУЛЬНЫЙ ЛИСТ
% ============================================================

\begin{titlepage}
\thispagestyle{empty}

\begin{center}

\vspace*{2.4cm}

{\Large\textcolor{accent}{\textbf{ЧЕК-ЛИСТ}}}

\vspace{0.9cm}

{\Huge\bfseries
Проценты
}

\vspace{0.35cm}

{\LARGE
Основные типы задач\\
и процентные изменения
}

\vspace{1.5cm}

{\large
7 класс
}

\vfill

{\large
\textbf{Лёвин Артём Александрович}\\[0.25em]
эксклюзивно для Кирилла Зиновьева
}

\vspace{0.7cm}

{\large \today}

\vfill

\begin{tikzpicture}[x=1cm,y=1cm]
    \draw[
        color=accent,
        line width=1.1pt
    ]
    (-6.7,0)
    .. controls (-4.7,0.65) and (-2.4,-0.55) ..
    (0,0)
    .. controls (2.4,0.55) and (4.7,-0.65) ..
    (6.7,0);
\end{tikzpicture}

\vspace{0.5cm}

\end{center}

\end{titlepage}

% ============================================================
% ОСНОВНОЙ ЧЕК-ЛИСТ
% ============================================================

\sectiontitle{1. Что показывает процент}

\key{Один процент} --- это одна сотая часть величины:

\[
1\%=\frac{1}{100}=0{,}01.
\]

Поэтому

\[
p\%=\frac{p}{100}.
\]

Например,

\[
30\%=\frac{30}{100}=0{,}3,
\qquad
3{,}2\%=\frac{3{,}2}{100}=0{,}032.
\]

Проценты удобно использовать для описания \key{доли}, \key{состава} и
\key{изменения величины}.

\sectiontitle{2. Три главных типа задач}

\subsectiontitle{Тип 1. Найти процент от числа}

Если требуется найти \(p\%\) от числа \(A\), используйте формулу

\[
\boxed{
A\cdot\frac{p}{100}
}
\]

или

\[
\boxed{
A\cdot 0{,}01p.
}
\]

\key{Пример.} Найдём \(25\%\) от \(36\):

\[
36\cdot\frac{25}{100}=9.
\]

Значит, \(25\%\) от \(36\) равны \(9\).

\subsectiontitle{Тип 2. Найти число по известному проценту}

Пусть \(p\%\) некоторого числа равны \(B\).

Тогда исходное число равно

\[
\boxed{
B:\frac{p}{100}
}
\]

или

\[
\boxed{
\frac{100B}{p}.
}
\]

\key{Пример.} \(30\%\) некоторого числа равны \(18\). Найдём это число:

\[
18:\frac{30}{100}
=
18\cdot\frac{100}{30}
=
60.
\]

Проверка:

\[
60\cdot\frac{30}{100}=18.
\]

\subsectiontitle{Тип 3. Найти новое значение после изменения на \(p\%\)}

Исходную величину удобно считать равной \(100\%\).

После увеличения на \(p\%\) новая величина составляет

\[
(100+p)\%.
\]

После уменьшения на \(p\%\) новая величина составляет

\[
(100-p)\%.
\]

Поэтому:

\[
\boxed{
A_{\text{нов}}
=
A_{\text{исх}}
\left(1+\frac{p}{100}\right)
}
\]

для увеличения и

\[
\boxed{
A_{\text{нов}}
=
A_{\text{исх}}
\left(1-\frac{p}{100}\right)
}
\]

для уменьшения.

\sectiontitle{3. Схема процентного изменения}

\begin{center}
\begin{tikzpicture}[
    node distance=2.3cm,
    >=Latex,
    every node/.style={
        font=\small,
        align=center
    }
]

\node[
    draw=accent,
    rounded corners=2pt,
    inner sep=7pt
] (base) {
    исходная величина\\
    \(100\%\)
};

\node[
    draw=accent,
    rounded corners=2pt,
    inner sep=7pt,
    above right=1.0cm and 2.5cm of base
] (up) {
    увеличение на \(p\%\)\\
    \((100+p)\%\)
};

\node[
    draw=accent,
    rounded corners=2pt,
    inner sep=7pt,
    below right=1.0cm and 2.5cm of base
] (down) {
    уменьшение на \(p\%\)\\
    \((100-p)\%\)
};

\draw[->,thick,color=darkaccent]
(base) -- node[above,sloped]{прибавляем \(p\%\)} (up);

\draw[->,thick,color=darkaccent]
(base) -- node[below,sloped]{вычитаем \(p\%\)} (down);

\end{tikzpicture}
\end{center}

\sectiontitle{4. Два способа решать задачи на увеличение}

Пусть первоначально собирали \(36\) центнеров пшеницы, затем урожай
увеличился на \(25\%\).

\subsectiontitle{Способ 1. Сначала найти величину увеличения}

\[
36\cdot\frac{25}{100}=9.
\]

Полученные \(9\) центнеров --- величина прироста.

Теперь найдём новый результат:

\[
36+9=45.
\]

\[
\boxed{45\text{ ц}}
\]

\subsectiontitle{Способ 2. Сразу перейти к новому проценту}

Исходное количество соответствует \(100\%\).

После увеличения на \(25\%\):

\[
100\%+25\%=125\%.
\]

Следовательно,

\[
36\cdot\frac{125}{100}=45.
\]

\[
\boxed{45\text{ ц}}
\]

Второй способ часто сокращает количество вычислений.

\sectiontitle{5. Скидка как уменьшение величины}

Книга стоила \(2000\) рублей. Скидка составляет \(20\%\).

После скидки остаётся

\[
100\%-20\%=80\%.
\]

Следовательно,

\[
2000\cdot\frac{80}{100}=1600.
\]

\[
\boxed{1600\text{ руб.}}
\]

То же вычисление можно записать через коэффициент:

\[
2000\cdot0{,}8=1600.
\]

\sectiontitle{6. Проценты и состав вещества}

Процент показывает, какую долю общей массы составляет некоторый компонент.

Если общая масса равна \(M\), а содержание компонента равно \(p\%\), то

\[
\boxed{
m_{\text{компонента}}
=
M\cdot\frac{p}{100}
}
\]

\subsectiontitle{Пример: состав молока}

В молоке содержится:

\begin{itemize}
    \item \(3{,}2\%\) жира;
    \item \(2{,}5\%\) белка;
    \item \(4{,}7\%\) углеводов.
\end{itemize}

Рассмотрим \(200\) г молока.

Жиры:

\[
200\cdot\frac{3{,}2}{100}=6{,}4\text{ г}.
\]

Белки:

\[
200\cdot\frac{2{,}5}{100}=5\text{ г}.
\]

Углеводы:

\[
200\cdot\frac{4{,}7}{100}=9{,}4\text{ г}.
\]

\begin{table}[h]
\centering
\caption{Состав \(200\) г молока}
\renewcommand{\arraystretch}{1.3}
\begin{tabularx}{0.76\linewidth}{Xcc}
\toprule
Компонент & Содержание & Масса \\
\midrule
Жиры & \(3{,}2\%\) & \(6{,}4\) г \\
Белки & \(2{,}5\%\) & \(5\) г \\
Углеводы & \(4{,}7\%\) & \(9{,}4\) г \\
\bottomrule
\end{tabularx}
\end{table}

Процентный состав сохраняется при любой выбранной общей массе.

Например, для \(500\) г того же молока получим:

\[
500\cdot\frac{3{,}2}{100}=16\text{ г},
\]

\[
500\cdot\frac{2{,}5}{100}=12{,}5\text{ г},
\]

\[
500\cdot\frac{4{,}7}{100}=23{,}5\text{ г}.
\]

\sectiontitle{7. Разность величин, заданных процентами}

За несколько книг заплатили \(320\) рублей.

Стоимость одной книги составляет \(30\%\) всей суммы, стоимость другой ---
\(45\%\).

\subsectiontitle{Способ 1. Найти обе стоимости}

Первая книга:

\[
320\cdot\frac{30}{100}=96\text{ руб.}
\]

Вторая книга:

\[
320\cdot\frac{45}{100}=144\text{ руб.}
\]

Разность:

\[
144-96=48\text{ руб.}
\]

\subsectiontitle{Способ 2. Сначала найти разность процентов}

Обе процентные величины относятся к одной общей сумме \(320\) рублей.

Поэтому

\[
45\%-30\%=15\%.
\]

Теперь найдём \(15\%\) от \(320\):

\[
320\cdot\frac{15}{100}=48.
\]

\[
\boxed{48\text{ руб.}}
\]

\key{Условие применимости.}
Складывать или вычитать процентные доли напрямую можно, когда они
относятся к одной и той же исходной величине.

\sectiontitle{8. Алгоритм выбора действия}

\begin{table}[h]
\centering
\caption{Как определить нужное действие}
\renewcommand{\arraystretch}{1.4}
\begin{tabularx}{\linewidth}{>{\bfseries}p{0.29\linewidth}X}
\toprule
Что требуется найти & Действие \\
\midrule

\(p\%\) от числа \(A\)
&
\[
A\cdot\frac{p}{100}
\]
\\

Число, \(p\%\) которого равны \(B\)
&
\[
B:\frac{p}{100}
\]
\\

Величину после увеличения на \(p\%\)
&
\[
A\left(1+\frac{p}{100}\right)
\]
\\

Величину после уменьшения на \(p\%\)
&
\[
A\left(1-\frac{p}{100}\right)
\]
\\

Разность двух долей \(p\%\) и \(q\%\) одной базы \(A\)
&
\[
A\cdot\frac{|p-q|}{100}
\]
\\

\bottomrule
\end{tabularx}
\end{table}

\sectiontitle{9. Полезные приёмы вычислений}

\begin{itemize}

    \item Перед умножением дробей сокращайте числители и знаменатели.

    \item Процент можно переводить в десятичную дробь:
    \[
    25\%=0{,}25,\qquad
    12{,}5\%=0{,}125.
    \]

    \item Для увеличения удобно сразу использовать коэффициент:
    \[
    +25\%
    \quad\Longrightarrow\quad
    \times1{,}25.
    \]

    \item Для уменьшения удобно сразу использовать коэффициент:
    \[
    -20\%
    \quad\Longrightarrow\quad
    \times0{,}8.
    \]

    \item После вычисления проверяйте единицы измерения:
    рубли, граммы, килограммы, центнеры и другие величины.

\end{itemize}

\sectiontitle{10. Типичные ошибки}

\begin{enumerate}

    \item \key{Ответом принимают только величину прироста.}

    При увеличении \(36\) на \(25\%\) число \(9\) показывает величину
    увеличения. Итоговое значение равно \(45\).

    \item \key{Путают два основных типа процентных задач.}

    Для поиска \(30\%\) от \(60\):

    \[
    60\cdot0{,}3.
    \]

    Для поиска числа, \(30\%\) которого равны \(18\):

    \[
    18:0{,}3.
    \]

    \item \key{Используют неверную исходную величину.}

    Всегда определяйте число, от которого считается процент.

    \item \key{Сравнивают проценты с разными базами как обычные числа.}

    Перед сложением или вычитанием процентных долей проверьте совпадение
    исходной величины.

    \item \key{Забывают смысл скидки.}

    Скидка \(20\%\) означает, что остаётся \(80\%\) исходной цены.

\end{enumerate}

\sectiontitle{11. Короткая самопроверка}

Перед записью ответа убедитесь:

\begin{itemize}
    \item определена исходная величина, соответствующая \(100\%\);
    \item понятно, требуется доля, исходное число или новое значение;
    \item процент переведён в дробь корректно;
    \item при увеличении учтено прибавление;
    \item при скидке учтено уменьшение;
    \item процентные доли складываются или вычитаются только при общей базе;
    \item итог снабжён нужными единицами измерения.
\end{itemize}

% ============================================================
% ТРЕНИРОВОЧНЫЙ БЛОК
% ============================================================

\clearpage
\sectiontitle{Тренировочный блок}

\textbf{Выполните задания самостоятельно.}
Старайтесь сначала определить тип процентной задачи и только затем
переходить к вычислениям.

\subsectiontitle{Средний уровень}

\begin{enumerate}

    \item Найдите \(18\%\) от числа \(250\).

    \item Найдите число, \(35\%\) которого равны \(84\).

    \item В \(350\) г молочного продукта содержится \(4\%\) жира.
    Сколько граммов жира содержится в продукте?

\end{enumerate}

\subsectiontitle{Выше среднего}

\begin{enumerate}[resume]

    \item Зарплата составляла \(48\,000\) рублей.
    Её увеличили на \(12{,}5\%\).
    Найдите новую зарплату.

    \item Книга стоила \(2400\) рублей.
    Магазин предоставил скидку \(15\%\).
    Найдите стоимость книги после скидки.

    \item За несколько товаров заплатили \(800\) рублей.
    Стоимость одного товара составляет \(27\%\) всей суммы,
    стоимость другого --- \(41\%\).
    На сколько рублей один товар дороже другого?

    \item Масса раствора равна \(600\) г.
    Содержание соли составляет \(8{,}5\%\).
    Найдите массу соли.

    \item После повышения цены на \(20\%\) товар стал стоить
    \(1560\) рублей.
    Найдите первоначальную цену.

    \item В прошлом году хозяйство собрало \(750\) центнеров зерна.
    В этом году количество уменьшилось на \(18\%\).
    Сколько центнеров зерна собрали в этом году?

\end{enumerate}

\subsectiontitle{Сложный уровень}

\begin{enumerate}[resume]

    \item Цена товара составляла \(4800\) рублей.
    Сначала её увеличили на \(25\%\), затем к полученной цене
    применили скидку \(20\%\).

    Найдите конечную цену товара и сравните её с первоначальной.

\end{enumerate}

% ============================================================
% ОТВЕТЫ
% ============================================================

\clearpage
\sectiontitle{Ответы к тренировочному блоку}

\begin{table}[h]
\centering
\caption{Ответы}
\renewcommand{\arraystretch}{1.45}
\begin{tabularx}{0.82\linewidth}{cX}
\toprule
\textbf{№} & \textbf{Ответ} \\
\midrule
1 & \(45\) \\
2 & \(240\) \\
3 & \(14\) г \\
4 & \(54\,000\) руб. \\
5 & \(2040\) руб. \\
6 & \(112\) руб. \\
7 & \(51\) г \\
8 & \(1300\) руб. \\
9 & \(615\) ц \\
10 & \(4800\) руб.; конечная цена совпадает с первоначальной \\
\bottomrule
\end{tabularx}
\end{table}

\vfill

\begin{center}
\textcolor{muted}{
\small
Главная идея: сначала определите величину, соответствующую \(100\%\),
затем выберите подходящую процентную модель.
}
\end{center}

\end{document}